\documentclass{ametsocV6.1}

%% AMS packages (ametsoc.sty is already loaded by the class)
\usepackage{graphicx}
\usepackage{amsmath,amssymb}
\usepackage{natbib}
\usepackage{caption}

% \journal{jcli}  % uncomment if using full ametsoc submission package

\begin{document}

\title{The relation of ENSO with global mean temperature}

\authors{Michael K. Tippett\aff{a}\correspondingauthor{M. K. Tippett,
    email@institution.edu} and Emily J. Becker\aff{b}}

\affiliation{
  \aff{a}{Department of Applied Physics and Applied
  Mathematics, Columbia University, New York, New York}\\
  \aff{b}{University of Miami Rosenstiel School for Marine, Earth,
  and Atmospheric Science, Miami, Florida}
}

\abstract{
  ENSO is a predictable contribution to global temperature. Previous studies found that the lag correlation between ENSO and
    global mean temperature peaks at 3--4 months with ENSO leading. Here we examined this relation with updated data and attention
    to seasonality. We found weaker relations than previously reported and that
    accounting for seasonality limits the peak correlations between
    ENSO and global mean temperature to January and February. The seasonality of the relation between ENSO and global mean
    temperature is explained by the seasonality of ENSO itself and by
    the seasonality of the atmospheric bridge mechanisms through which
    ENSO-related circulation drives remote SST changes.
}

\maketitle

\section{Introduction}

People like to talk about global mean temperature will do in
    the future. \citet{Tippett2024trends} found that the skill of initialized
    seasonal climate forecasts in predicting global mean temperature
    was primarily attributable to ENSO. When ENSO was regressed out
    the forecasts, the remaining skill was little more than that of a
    linear trend, which suggests that ENSO is the primary source of
    global mean temperature predictability in current seasonal
    forecast systems. Since ENSO can be skillfully predicted a few seasons ahead,
    better understanding of the influence of ENSO on global mean
    temperature may improve our ability to anticipate short-term
    variations in global mean temperature beyond the upward trend due
    to anthropogenic forcing.

A notable feature of the ENSO–global-temperature relationship
    is the delayed response of global temperature response to ENSO. In an early study, \citet{Angell1981comparison} compared
    variations of an equatorial eastern Pacific SST index
    (0--10\textdegree S, 180\textdegree--90\textdegree W) with
    variations in atmospheric temperature from a 63-station global
    radiosonde network. They found that tropical tropospheric
    temperature (0--16km average) as well as surface temperature was
    maximally correlated with the equatorial Pacific SST index at
    about a two-season lag.  Relations with atmospheric temperature in
    temperate latitudes were weaker.

  

\citet{Angell2000tropospheric} used this lagged relationship
    to account for variability in global temperatures due to ENSO and
    thus improve estimates of temperature trends. \citet{Lean2008natural} took a similar approach to separating
    forced and natural variability in surface temperature. They
    included an ENSO index, lagged by four months, in a multiple linear
    regression fit to global mean temperature that also included
    volcanic aerosols, solar irradiance, and anthropogenic
    forcing. This framework allowed them to attribute about
    0.23\textdegree C of global mean temperature warming to the 1997
    ``super'' El Ni\~no, approximately 0.096\textdegree C of global
    mean temperature for each degree change in Nino-3.4. Similarly, \citet{Foster2011global} fit global mean
    temperature by a multiple linear regression that included a
    standardized ENSO index, aerosol optical depth, total solar
    irradiance, and a linear time trend over the period
    1979--2010. Depending on temperature dataset, global mean
    temperature lagged their ENSO index by 2--5 months, and global
    mean temperature changed approximately 0.075 \textdegree C per
    standardized ENSO index unit.

  
  

\citet{Trenberth2002evolution} examined spatially-resolved
    diabatic processes in the relationship between ENSO and global
    temperature. They concluded that delayed surface warming outside of the
    tropical Pacific was due to persistent atmospheric circulation
    changes forced from the tropical Pacific. They found that the lag correlation between Nino-3.4 and
    global mean surface temperature peaked at three to four months,
    depending on period. They speculated that differing ENSO indices
    might be a reason for the longer two season lag priviously found
    by \citet{Angell1981comparison}. They computed a three-month lagged regression coefficient of
    0.094\textdegree C per unit of standardized Nino-3.4 for the
    period 1950--1998. \citet{Kumar2003nature} focused on the delayed
    \textit{atmospheric} response to El Ni\~no and explained the delay
    as being due to a delayed tropical rainfall response to ENSO which
    in turn was attributed to the seasonality of SST variability in the
    tropical eastern Pacific. \citet{Su2005mechanisms} questioned this explanation and found
    that the lagged response of tropospheric temperature in models to
    ENSO forcing depended on mixed layer depth and showed little lag
    sensitivity to the seasonal cycle of SST.

  
 
  

The ``atmospheric bridge'' perspective introduced in the 1990s
    provides a physical explanation for the delayed remote response to
    ENSO \citep[see][ for a review]{Alexander2002atmospheric}. The atmospheric bridge refers to the set of teleconnections
    and mechanisms through which ENSO-related SST anomalies in the
    tropical Pacific force SST anomalies in remote ocean basins via
    changes in radiative and heat fluxes. The atmospheric bridge responses in the tropics (Indian Ocean,
    South China Sea, Atlantic) have the same sign as the tropical
    Pacific ENSO signal and thus contribute constructively to global
    mean temperature anomalies. In the atmospheric bridge, the delayed remote response is due
    to the rapid (on the order of two weeks) atmospheric adjustment to
    tropical Pacific SST anomalies being followed by the slower
    integration of the resulting flux anomalies by the ocean mixed
    layer \citep{Klein1999remote}. Seasonality of ENSO and its
    teleconnections create seasonality in the remote responses.

There is an inconsistency between studies which pooled all
    months together to examine the lagged relationship of ENSO and
    global temperature \citep[e.g.,][]{Trenberth2002evolution,
      Lean2008natural, Foster2011global}, thus ignoring seasonality,
    and atmospheric bridge studies which highlight seasonality by
    focusing on the ENSO peak and the delayed remote response. In addition, the most recent spatially resolved analysis of the
    relationship between ENSO and global temperature used data ending
    in 1998 or so. Time for an update!

\section{Data and methods}

Our analysis used monthly data over the period 1979--2025. The Ni\~no-3.4 index is the average of sea surface temperature
    (SST) in the tropical Pacific box  170\textdegree W--120\textdegree W,
    5\textdegree S--5\textdegree N \citep{Barnston1997} and was
    computed from ERSSTv5 \citep{Huang:ERSST5}. Warm and cool ENSO years were identified from the NOAA Oceanic
    Nino Index (ONI, version 5) based on December--February Nino-3.4
    conditions. For each event, composites were formed using a
    16-month window beginning in June (year 0) before the boreal
    winter peak and extending through September of the following year
    (year +1). Thus, for example, the 1997/98 El Ni\~no corresponds to
    the period June 1997 through September 1998. Warm ENSO years were
    1979, 1982, 1986, 1987, 1991, 1994, 1997, 2002, 2004, 2006, 2009,
    2014, 2015, 2018, 2019, and 2023. Cool ENSO years were 1983, 1984,
    1988, 1995, 1998, 1999, 2000, 2005, 2007, 2008, 2010, 2011, 2017,
    2020, 2021, and 2022.

  

Averages of near-surface temperature were computed from the
    NOAA global surface temperature version 6.1
    \citep[NOAAGlobalTemp;][]{Huang2022improvements}. 5$\times$5
    resolution which means the tropics are $\pm 22.5$.

  

Box definitions. Boxes on map are in Fig.\ \ref{fig:map-box}.

  

Fluxes etc.\ are taken from the fifth generation ECMWF
    atmospheric reanalysis of the global climate
    \citep[ERA5;][]{Hersbach2020era5}.

Anomalies and detrending details. Tendencies were computed by first differencing in time. Each
    tendency was labeled by the later of the two months used in the
    difference, so that the tendency value at month $t$ represents the
    change from month $t-1$ to month $t$. Accordingly, the first month
    has no defined tendency value.

  

Lag correlations $C(L)$ were computed as the correlation
    between $Y_{t+L}$ and $X_t$, where $L$ denotes lag in
    months. Here, $X_t$ and $Y_t$ represent, for example, the Niño-3.4
    index and global mean temperature at month $t$. Lags from $0$ to
    $18$ months were considered, and all calendar months were included
    in the calculation. Month-dependent lag correlations $C(M,L)$ were computed
    separately for each calendar month $M$ ($M=1,2,\ldots,12$) as the
    correlation between $Y_{t+L}$ and $X_t$ for times $t$ in calendar
    month $M$. These month-dependent lag correlations were plotted as
    a function of the target calendar month $M+L$. Lag regression coefficients were computed in the same way,
    replacing correlation by regression.

  

The statistical significance of correlations, and equivalently
    of regression coefficients, was assessed using a block bootstrap
    with 12-month resampling segments (1000 samples with replacement)
    to account for serial correlation. For tendencies, standard
    bootstrap resampling was used because serial correlation was not
    evident. We used applied a permutation test to assess the statistical
    significance of ENSO composites.

\section{Results}

\textbf{Overview}
  

Lag correlations of Nino-3.4 with global mean temperature peak
    at three months and are weaker than previously reported. Removing
    trends increases the strength of correlations. Nino-3.4 is
    significantly related with global mean temperature tendencies only
    at zero lag, consistent with the atmospheric bridge mechanism.
    (Fig.\ \ref{fig:no-seasonality}). Month-dependent lag correlations of Nino-3.4 with global mean
    temperature peak in January and are statistically insignificant in
    most other months.  Month-dependent lag correlations of Nino-3.4
    with tropical mean temperature peak in February and are
    statistically significant across the ENSO cycle.
    (Fig. \ref{fig:with-seasonality}). ENSO composites show that in the zonal mean temperature
    anomalies are mostly confined to the tropics. Within the tropics,
    ENSO-related temperature anomalies outside the tropical Pacific
    emerge at different times in the ENSO cycle, consistent with
    atmospheric bridge mechanisms. (Fig. \ref{fig:composite}). Area-weighted composites of tropical basin averages explain
    the peak timings. (Figs.\ \ref{fig:basin-composite} and
    \ref{fig:basin-corr}). Radiative and heat flux composites (primarily solar and
    latent) do their thing. (Fig.\
    \ref{fig:flux_composites_solar_latent}).

\begin{figure}[h]
  \centering
  \includegraphics[width=\linewidth]{/Users/tippett/notebooks/NMME/enso-t2m-figures/lag-corr-reg-without-seasonality-mini.pdf}
  \caption{Lag (a) correlation and (b)
      regression coefficients of Nino-3.4 with global and tropical
      mean temperature.  Lag (c) correlation and (d) regression
      coefficients of Nino-3.4 with global and tropical mean
      temperature tendencies.  Filled circles indicate statistical
      significance at the 5\% level.  Results for detrended time
      series are shown in dark colors.}
  \label{fig:no-seasonality}
\end{figure}

  
   
    

Computed over all months, Nino-3.4 lag correlation and regression
      coefficients peak at a lag of three months for global mean
      temperature and at a lag of two months for tropical mean
      temperature (Figs.\ \ref{fig:no-seasonality}a, b). Correlations of Nino-3.4 with global mean temperature are
      substantially lower than those in \citet[][; their Fig.\
      2]{Trenberth2002evolution} and are statistically significant
      only at a lag of three months (5\% level). Statistically significant lag correlations of Nino-3.4 with
      tropical mean temperature extend to 10 months. Detrending the data increases the strength of the
      correlations substantially. The global mean temperature
      correlation with Nino-3.4 at lag three increases from 0.21 to
      0.48 after detrending. The regression coefficients change little
      with detrending because the long-term Ni\~no-3.4 trend is small. On average, a one degree increase in Nino-3.4 is associated
      with increases of 0.07\textdegree C and 0.16\textdegree C at
      lags of two to four months in detrended global and tropical mean
      temperatures, respectively. Nino-3.4 is significantly related only at zero lag with
      global mean temperature tendencies and only at lags zero and one
      with tropical mean temperature tendencies (Figs.\
      \ref{fig:no-seasonality}c, d). This behavior is consistent with the atmospheric bridge
      mechanism \citep{Klein1999remote} whereby ENSO-forced
      circulation anomalies drive changes in SST with little delay. Detrending the data does not visibly change tendency
      correlation and regression coefficients because differencing the
      data effectively detrends it (Figs.\ \ref{fig:no-seasonality}c,
      d). See Figs.\ \ref{fig:no-seasonality-cor-map} and
      \ref{fig:no-seasonality-cor-map-diff} for lag correlation maps.

\begin{figure}[h]
  \centering
  \includegraphics[width=\linewidth]{/Users/tippett/notebooks/NMME/enso-t2m-figures/lag-corr-reg-with-seasonality-2-lines.pdf}
  \caption{Month-dependent lag
      correlations and regression coefficients of Nino-3.4 with (a, b)
      tropical mean temperature and (c, d) global mean
      temperature. The x-axis shows the target month (start month +
      lag), and each line corresponds to a different start month
      May--April extending 18 months. Gray (black) dots indicate
      statistical significance at the 5\% (1\%)
      level.}
  \label{fig:with-seasonality}
\end{figure}

  
   
    

Month-dependent lag correlation and regression coefficients
      of Nino-3.4 with tropical and global mean temperature are phase
      locked to the annual cycle (Fig.\ \ref{fig:with-seasonality}). They depend on target month primarily rather than lag. For start months after May, tropical mean temperature
      correlation and regression coefficients peak in February, and
      those for global mean temperature peak in January. Peak regression coefficient values exceed the all-months
      ones in Fig.\ \ref{fig:no-seasonality} by roughly a factor of
      two. Statistically significant lag correlations between Nino-3.4
      and global mean temperature are mostly limited to the target months
      of January and February. Statistically significant lag correlations between Nino-3.4
      and tropical mean extend to the end of the ENSO cycle in
      May/June. For a given target month, the lag regression coefficient
      varies more with start month than the lag correlation
      coefficients do because the lag regression coefficients depend
      inversely on the Nino-3.4 start month variance, which has a
      strong seasonal cycle. Detrending increases correlation and statistical
      significance (Fig.\ \ref{fig:corr-with-seasonality-chiclet}). The average of the month-dependent lag correlations over
      start month approximately recovers the the all-months lag
      correlations (Fig.\ \ref{fig:pooled-vs-ave}).

\begin{figure}[h]
  \centering
  \includegraphics[width=0.49\linewidth]{/Users/tippett/notebooks/NMME/enso-t2m-figures/zonal-mean-composite-window.pdf}
  \includegraphics[width=0.49\linewidth]{/Users/tippett/notebooks/NMME/enso-t2m-figures/mean-tropics-composite-window.pdf}
  \caption{Composites of (a, c) zonal and (b, d) tropical mean temperature
      anomaly from June (year 0) through September (year +1) during
      warm and (b) cool ENSO years. White (black) dots indicate
      statistical significance at the 5\% (1\%) level. Dashed lines in
      panels b and d show Indian (green), South China Sea (purple),
      tropical Pacific (red), South America (orange), tropical
      Atlantic (blue), and Africa (olive)
      sectors.}
  \label{fig:composite}
\end{figure}

To understand the spatially-resolved sources of seasonality in
    the month-dependent relations, we computed warm and cool ENSO
    composites of zonal mean (Figs.\ \ref{fig:composite}a, c) and
    tropical mean (Fig.\ \ref{fig:composite}b, c) temperature
    anomalies over a 16-month window spanning the ENSO development and
    decay life cycle (June year 0 preceding the boreal winter peak
    through September of the following year +1). Statistically significant zonal mean anomalies are mostly
    confined to the tropics. Initially in the ENSO cycle, zonal mean anomalies are confined
    to the equation. They extend northward briefly in October and then
    again in February. Temperature anomalies extend southward from the
    equator in November where they persist into the following boreal
    summer. Equatorial anomalies weaken and fade by May/June following the
    boreal winter peak as the ENSO cycle ends. Off-equatorial tropical
    anomalies persist longer, especially in the warm ENSO composites.

  

The tropical mean composites decompose the zonal mean signals
    by longitude. The dominant ENSO signal in terms of duration and
    spatial extent is the direct one in the central-eastern tropical
    Pacific where it persists for the full ENSO cycle. Significant anomalies in the Indian sector appear in late
    boreal fall (October/November) and persist for nearly a year. South China Sea sector anomalies show warm ENSO composite
    phase signals lasting from December to September. South American sector composite anomalies are widespread by
    November and last until the following June, with warm phase
    anomalies lasting longer. Temperature anomalies in Atlantic sectors also show
    substantial differences between warm and cool phases, with the
    warm-phase signal having broad longitudinal extent in February
    that persists several months while the cool-phase signal mostly
    statistically insignificant. Warm-ENSO anomalies in in the tropical Atlantic expand
    eastward into Africa in February, together with anomalies in the
    Indian and South China Sea sectors that have their maximum
    longitudinal extent in February. The generally more muted (duration and extent) cool ENSO
    composites may reflect imperfect removal of warming trends.

\begin{figure}[h]
  \centering
  \includegraphics[width=\linewidth]{/Users/tippett/notebooks/NMME/enso-t2m-figures/basin-mean-tropics-composite-window.pdf}
  \caption{Area-weighted basin composites
      of temperature.}
  \label{fig:basin-composite}
\end{figure}
 
    

Warm ENSO tropical mean temperature composites peak in
      February while cool ones show a less well defined peak in
      January--February. Temperatures outside the tropical Pacific are responsible
      for this timing. All the remote sector warm ENSO temperature composites peak
      in February except the South America one.

\begin{figure}[h]
  \centering
  \includegraphics[width=\linewidth]{/Users/tippett/notebooks/NMME/enso-t2m-figures/lag-cor-with-seasonality-2-lines-regional.pdf}
  \caption{Lag correlations of Nino-3.4
      with basin-averaged temperature.}
  \label{fig:basin-corr}
\end{figure}
 
    

Lagged correlations give the same story as the composites. Peak in tropical Atlantic and Africa averages. Regression versions in Fig.\ \ref{fig:basin-reg}. Chiclet versions in Figs.\ \ref{fig:basin-corr-2} and
      \ref{fig:basin-reg-2}.

\begin{figure}[h]
  \centering
  \includegraphics[width=0.49\linewidth]{/Users/tippett/notebooks/NMME/enso-t2m-figures/mean-tropics-composite-window-ssr.pdf}
  \includegraphics[width=0.49\linewidth]{/Users/tippett/notebooks/NMME/enso-t2m-figures/mean-tropics-composite-window-slhf.pdf}
  \caption{Composites of tropical mean (a, c) surface net solar radiation
      and (b, d) surface latent heat flux anomaly from June (year 0)
      through September (year +1) during warm and (b) cool ENSO
      years. (Coarsened to five degree in longitude). White (black)
      dots indicate statistical significance at the 5\% (1\%)
      level. Dashed lines in panels b and d show Indian (green), South
      China Seas (purple), tropical Pacific (red), South America
      (orange), tropical Atlantic (blue), and Africa (olive)
      sectors.}
  \label{fig:flux_composites_solar_latent}
\end{figure}

Figure~\ref{fig:flux_composites_solar_latent} shows composites
    of surface net solar radiation (panels a, c) and surface latent
    heat flux (panels b, d), which are the dominant individual
    components of the net radiative and net heat flux composites,
    respectively (Fig.\ \ref{fig:flux_composites}). Latent heat flux is the dominant component of the net
    turbulent heat flux (sensible heat flux being secondary over
    tropical oceans due to the relatively small tropical air–sea
    temperature contrast). The warm and cool ENSO composites are approximately
    antisymmetric, consistent with a roughly linear atmospheric
    response to ENSO-related SST forcing. During warm ENSO, net solar radiation in the tropical Pacific
    is significantly reduced east of the date line during the peak of
    the ENSO cycle, coinciding with enhanced deep convection and
    increased cloud cover over anomalously warm SSTs and a shifted
    rising branch of the Walker circulation (see Fig.\
    \ref{fig:tcc_omega500_composites} for ENSO composites of cloud
    cover and omega 500). The descending circulation to the west
    accompanies increases in solar radiation. West of 150E solar radiation anomalies are significantly
    positive September--March together with reduced cloud cover and
    enhanced ascending motion. This confirms the cloud--shortwave
    mechanism as the radiative pathway of the atmospheric bridge to
    the eastern Indian Ocean and South China Sea
    \citep{Klein1999remote}. Surface latent heating in the tropical Pacific around 120W is
    mostly a negative feedback on ENSO SST anomalies due to increased
    evaporative heat loss from the warmer ocean surface, which
    counteracts weakened trades.

    In the eastern Indian Ocean and South China Sea, weakened trades
    during warm ENSO conditions reduce evaporation and increase latent
    heating but are limited to November and December. They weaker and
    less extensively significant than the solar radiation anomalies in
    the same sector. This indicates that the radiative
    (cloud--shortwave) pathway is the larger contributor to Indian
    Ocean warming, with the wind--evaporation pathway playing a
    supporting role. Over the tropical Atlantic latent heat flux anomalies are
    positive during boreal December through February, with SST
    expected to follow at a lag. There is essentially no solar
    radiation signal over the Atlantic, indicating that the
    atmospheric bridge to the tropical Atlantic operates primarily
    through the evaporative pathway rather than the cloud--radiative
    pathway. The warm and cool surface latent heating composites are
    similar over the Atlantic even though the temperature ones in
    Fig.\ \ref{fig:composite}are not. Maps of the correlation of Nino-3.4 with surface net solar
    radiation and surface latent heat flux by month to confirm
    seasonality maps are shown in Figs.\ \ref{fig:bridge-ind} and
    \ref{fig:bridge-atl}.

\section{Conclusions and discussion}

In early 2026, the WMO noted ``The most recent El Ni\~no, in
    2023-24, was one of the five strongest on record and it played a
    role in the record global temperatures we saw in 2024''
    \citep{WMO2026enso}.

NOAA Extended Reconstructed SST V5 (ERSSTv5) data are provided
  by the NOAA PSL, Boulder, Colorado, USA, from their website
  \url{https://psl.noaa.gov}. NOAAGlobalTemp version 6.1 data are
  from \url{https://www.ncei.noaa.gov/data/noaa-global-surface-temperature/v6.1/access/gridded/NOAAGlobalTemp_v6.1.0_gridded_s185001_e202602_c20260308T114550.nc}. ERA5 data are available from the Climate Data Store
  \url{https://cds.climate.copernicus.eu/}.

\bibliographystyle{ametsocV6}
  \bibliography{all}

\section{Supplemental/not shown}

\renewcommand\thefigure{S\arabic{figure}}

\setcounter{figure}{0}

\begin{figure}
  \includegraphics[width=\linewidth]{/Users/tippett/notebooks/NMME/enso-t2m-figures/map-of-boxes.pdf}
  \caption{}\label{fig:map-box}
\end{figure}

\begin{figure}[h]
  \centering
  \includegraphics[width=\linewidth]{/Users/tippett/notebooks/NMME/enso-t2m-figures/lag-corr-maps.png}
  \caption{Lag correlation of Nino-3.4 with 2m
      temperature. White (black) dots indicate statistical
      significance at the 5\% (1\%) level. The global mean of the
      correlation is shown in the title with bold indicating the lag
      with maximum value for that start
      month.}
  \label{fig:no-seasonality-cor-map}
\end{figure}

\begin{figure}[h]
  \centering
  \includegraphics[width=\linewidth]{/Users/tippett/notebooks/NMME/enso-t2m-figures/lag-corr-maps-diff.png}
  \caption{Lag correlation of Nino-3.4 with
      temperature \textbf{tendencies}. White (black) dots indicate
      statistical significance at the 5\% (1\%) level. The global mean
      of the correlation is shown in the title with bold indicating
      the lag with maximum value.}
  \label{fig:no-seasonality-cor-map-diff}
\end{figure}

\begin{figure}[h]
  \centering
  \includegraphics[width=\linewidth]{/Users/tippett/notebooks/NMME/enso-t2m-figures/lag-corr-with-seasonality-2-target.pdf}
  \caption{Lag correlations as a function of
      start month and target month of Nino-3.4 with (a) tropical mean
      temperature, (b) detrended tropical mean temperature, (c) global
      mean temperature, (d) detrended global mean temperature. White
      (black) dots indicate statistical significance at the 5\% (1\%)
      level. Detrended analysis detrends both time
      series.}
  \label{fig:corr-with-seasonality-chiclet}
\end{figure}

\begin{figure}[h]
  \centering
  \includegraphics[width=\linewidth]{/Users/tippett/notebooks/NMME/enso-t2m-figures/lag-corr-with-ave-seasonality-mini.pdf}
  \caption{All-months lag correlations and
      averages of month-dependent lag correlations of Nino-3.4 with
      (a) tropical mean temperature and (b) global mean
      temperature. Results for detrended time series are shown in dark
      colors.}
  \label{fig:pooled-vs-ave}
\end{figure}

  
   
    

The average of the month-dependent lag correlations matches
      well the all-month lag correlations (Fig.\
      \ref{fig:pooled-vs-ave}). Math says it doesn't have to work out so nicely but it
      happens to. Previously reported peaks at lag-3 for global mean
      temperature are only correct in an average (sloppy) sense that
      obscures what is really going on.

\begin{figure}[h]
  \centering
  \includegraphics[width=\linewidth]{/Users/tippett/notebooks/NMME/enso-t2m-figures/lag-reg-with-seasonality-2-lines-regional.pdf}
  \caption{Lag regression of Nino-3.4
      with basin averaged temperature.}
  \label{fig:basin-reg}
\end{figure}

\begin{figure}[h]
  \centering
  \includegraphics[width=\linewidth]{/Users/tippett/notebooks/NMME/enso-t2m-figures/lag-corr-with-seasonality-2-target-regional.pdf}
  \caption{Lag correlations of Nino-3.4
      with basin averaged temperature.}
  \label{fig:basin-corr-2}
\end{figure}

\begin{figure}[h]
  \centering
  \includegraphics[width=\linewidth]{/Users/tippett/notebooks/NMME/enso-t2m-figures/lag-reg-with-seasonality-2-target-regional.pdf}
  \caption{Lag regressions of Nino-3.4
      with basin averaged temperature.}
  \label{fig:basin-reg-2}
\end{figure}

\begin{figure}[h]
  \centering
  \includegraphics[width=0.49\linewidth]{/Users/tippett/notebooks/NMME/enso-t2m-figures/mean-tropics-composite-window-sr.pdf}
  \includegraphics[width=0.49\linewidth]{/Users/tippett/notebooks/NMME/enso-t2m-figures/mean-tropics-composite-window-shf.pdf}
  \caption{Composites of tropical mean (a, c) net surface radiation
      (shortwave plus longwave) and (b, d) surface heat flux (latent
      plus sensible) anomaly from June (year 0) through September
      (year +1) during warm and (b) cool ENSO years. (Coarsened to
      five degree in longitude). White (black) dots indicate
      statistical significance at the 5\% (1\%) level. Dashed lines in
      panels b and d show Indian (green), South China Seas (purple),
      tropical Pacific (red), South America (orange), tropical
      Atlantic (blue), and Africa (olive)
      sectors.}
  \label{fig:flux_composites}
\end{figure}

\begin{figure}[h]
  \centering
  \includegraphics[width=0.49\linewidth]{/Users/tippett/notebooks/NMME/enso-t2m-figures/mean-tropics-composite-window-tcc.pdf}
  \includegraphics[width=0.49\linewidth]{/Users/tippett/notebooks/NMME/enso-t2m-figures/mean-tropics-composite-window-omega500.pdf}
  \caption{Composites of tropical mean (a, c) total cloud cover and (b, d)
      omega 500 from June (year 0) through
      September (year +1) during warm and (b) cool ENSO
      years. (Coarsened to five degree in longitude). White (black)
      dots indicate statistical significance at the 5\% (1\%)
      level. Dashed lines in panels b and d show Indian (green), South
      China Seas (purple), tropical Pacific (red), South America
      (orange), tropical Atlantic (blue), and Africa (olive)
      sectors.}
  \label{fig:tcc_omega500_composites}
\end{figure}

\begin{figure}[h]
  \centering
  \includegraphics[width=0.49\linewidth]{/Users/tippett/notebooks/NMME/enso-t2m-figures/map-surface-net-solar-radiation-n34-corr-indo.png}
  \includegraphics[width=0.49\linewidth]{/Users/tippett/notebooks/NMME/enso-t2m-figures/map-surface-latent-heat-flux-n34-corr-indo.png}
  \caption{}
  \label{fig:bridge-ind}
\end{figure}

\begin{figure}[h]
  \centering
  \includegraphics[width=0.49\linewidth]{/Users/tippett/notebooks/NMME/enso-t2m-figures/map-surface-net-solar-radiation-n34-corr-atl.png}
  \includegraphics[width=0.49\linewidth]{/Users/tippett/notebooks/NMME/enso-t2m-figures/map-surface-latent-heat-flux-n34-corr-atl.png}
  \caption{}
  \label{fig:bridge-atl}
\end{figure}

\end{document}
